%-------------------------
% Resume in Latex
% Author : Jake Gutierrez
% Based off of: https://github.com/sb2nov/resume
% License : MIT
%------------------------

\documentclass[letterpaper,11pt]{article}
\usepackage{fancyhdr}
\usepackage{latexsym}
\usepackage{cite}
\usepackage[empty]{fullpage}
\usepackage{titlesec}
\usepackage{marvosym}
\usepackage[usenames,dvipsnames]{color}
\usepackage{verbatim}
\usepackage{enumitem}
\usepackage[hidelinks]{hyperref}
\usepackage{fancyhdr}
\usepackage[english]{babel}
\usepackage{tabularx}
\usepackage{fontawesome5}
\usepackage{multicol}
\setlength{\multicolsep}{-3.0pt}
\setlength{\columnsep}{-1pt}
\input{glyphtounicode}
\usepackage{amsmath}


%----------FONT OPTIONS----------
% sans-serif
% \usepackage[sfdefault]{FiraSans}
% \usepackage[sfdefault]{roboto}
% \usepackage[sfdefault]{noto-sans}
% \usepackage[default]{sourcesanspro}

% serif
% \usepackage{CormorantGaramond}
% \usepackage{charter}


\pagestyle{fancy}
\fancyhf{} % clear all header and footer fields
\fancyhead[L]{}
\fancyfoot{}
\renewcommand{\headrulewidth}{0pt}
\renewcommand{\footrulewidth}{0pt}

% Adjust margins
\addtolength{\oddsidemargin}{-0.6in}
\addtolength{\evensidemargin}{-0.5in}
\addtolength{\textwidth}{1.19in}
\addtolength{\topmargin}{-.7in}
\addtolength{\textheight}{1.4in}

\urlstyle{same}

\raggedbottom
\raggedright
\setlength{\tabcolsep}{0in}

% Sections formatting
\titleformat{\section}{
  \vspace{-4pt}\scshape\raggedright\large\bfseries
}{}{0em}{}[\color{black}\titlerule \vspace{-5pt}]

% Ensure that generate pdf is machine readable/ATS parsable
\pdfgentounicode=1

%-------------------------
% Custom commands
% new command for publication
\newcommand{\publicationentry}[3]{%
    \noindent
    {#1,} #2. \emph{#3}\vspace{-5pt}
}

\newcommand{\resumeItem}[1]{
  \item\small{
    {#1 \vspace{-2pt}}
  }
}

\newcommand{\classesList}[4]{
    \item\small{
        {#1 #2 #3 #4 \vspace{-2pt}}
  }
}

\newcommand{\resumeSubheading}[4]{
  \vspace{-2pt}\item\vspace{2pt}
    \begin{tabular*}{1.0\textwidth}[t]{l@{\extracolsep{\fill}}r}
      \textbf{#1} & \textbf{\small #2} \\
      \textit{\small#3} & \textit{\small #4} \\
    
    \end{tabular*}\vspace{-5pt}
}

\newcommand{\resumeSubSubheading}[2]{
    \item
    \begin{tabular*}{0.97\textwidth}{l@{\extracolsep{\fill}}r}
      \textit{\small#1} & \textit{\small #2} \\
    \end{tabular*}\vspace{-7pt}
}

\newcommand{\resumeProjectHeading}[2]{
    \item
    \begin{tabular*}{1.001\textwidth}{l@{\extracolsep{\fill}}r}
      \small#1 & \textbf{\small #2}\\
    \end{tabular*}\vspace{-7pt}
}

\newcommand{\resumeSubItem}[1]{\resumeItem{#1}\vspace{-4pt}}

\renewcommand\labelitemi{$\vcenter{\hbox{\tiny$\bullet$}}$}
\renewcommand\labelitemii{$\vcenter{\hbox{\tiny$\bullet$}}$}

\newcommand{\resumeSubHeadingListStart}{\begin{itemize}[leftmargin=0.0in, label={}]}
\newcommand{\resumeSubHeadingListEnd}{\end{itemize}}
\newcommand{\resumeItemListStart}{\begin{itemize}}
\newcommand{\resumeItemListEnd}{\end{itemize}\vspace{-5pt}}

%-------------------------------------------
%%%%%%  RESUME STARTS HERE  %%%%%%%%%%%%%%%%%%%%%%%%%%%%


\begin{document}

%----------HEADING----------
% \begin{tabular*}{\textwidth}{l@{\extracolsep{\fill}}r}
%   \textbf{\href{http://sourabhbajaj.com/}{\Large Sourabh Bajaj}} & Email : \href{mailto:sourabh@sourabhbajaj.com}{sourabh@sourabhbajaj.com}\\
%   \href{http://sourabhbajaj.com/}{http://www.sourabhbajaj.com} & Mobile : +1-123-456-7890 \\
% \end{tabular*}

\begin{center}
    {\Huge \scshape Zheng Zheng} \\ \vspace{1pt}
    Arlington, TX, USA, 76019 \\ \vspace{1pt}
    \small \raisebox{-0.1\height}\faPhone\ 401-215-4595 ~ \href{mailto:zxz7934@mavs.uta.edu}{\raisebox{-0.2\height}\faEnvelope\ \underline{zxz7934@mavs.uta.edu}} ~ 
    \href{https://www.linkedin.com/in/zheng-zheng-082420217/}{\raisebox{-0.2\height}\faLinkedin\ \underline{linkedin}}  ~
    \href{https://zz9tf.github.io/}{\raisebox{-0.2\height}\faGithub\ \underline{portfolio}}
    \vspace{-7pt}
\end{center}


%-----------EDUCATION-----------
\section{Education}
  \resumeSubHeadingListStart

    \resumeSubheading
      {University of Texas at Arlington}{Aug. 2024 -- Present}
      {Ph.D. in Computer Science}{Advisor: Prof. Junzhou Huang}
      \resumeItemListStart
        \resumeItem{Developed multimodal alignment framework for 7 modalities under missing-modality settings in pathology prediction.}
            
        \resumeItem{Analyzed and designed algorithms for multi-type gradient conflict in multi-task learning, including structural subspace conflict, batch-induced stochastic noise conflict, and dominant gradient imbalance.}

        \resumeItem{Designed and fine-tuned protein foundation models (PTM) and implemented structured knowledge injection mechanisms to enhance representation adaptation and transferability.}
      \resumeItemListEnd
      
    \vspace{2pt}

    \resumeSubheading
      {Brandeis University}{Sep. 2021 -- Dec. 2023}
      {Master of Science in Computer Science}{}
      \resumeItemListStart
      %     \resumeItem{Relevant Courses: Deep Learning (A), Fundamentals of Natural Language Processing (A), Data Structures and the Fundamentals of Computing (A), Advanced Programming Techniques in Java (A-), Discrete Structures (A), Theory of Computation (A), 3-D Animation(A), Operating Systems (B)}
      \resumeItemListEnd
      
    \vspace{2pt}
    
    \resumeSubheading
      {Missouri University of Science and Technology }{Sep. 2019 -- May 2021}
      {Bachelor of Science in Geology and Geophysics }{}
      \resumeItemListStart
      %     \resumeItem{Relevant Courses: Computational Geophysics (A), Discrete Math (A), Global Tectonics (A)}
      \resumeItemListEnd
      

  \resumeSubHeadingListEnd
  \vspace{-7pt}

%-----------Publication-----------
% \section{Publication}
%     \begingroup
%     \renewcommand{\section}[2]{}%

%     \begin{thebibliography}{9}

%     \bibitem{Zhang2023}
%     Pengtao Zhang, \textbf{Zheng Zheng}, and Junlin Zhang. FiBiNet++: Reducing model size by low rank feature interaction layer for CTR prediction. In \textit{Proceedings of the 32nd ACM International Conference on Information and Knowledge Management}, \href{https://dl.acm.org/doi/10.1145/3583780.3615242}{\color{blue}CIKM' 23}, pages 4425–4429, New York, NY, USA, 2023. Association for Computing Machinery.
    
%     \end{thebibliography}
%     \endgroup
\section{Publications}
    \begingroup
    \renewcommand{\section}[2]{}%

    \begin{thebibliography}{9}

    \bibitem{Zheng2026}
    \textbf{Zheng Zheng}, Y. Guo, X. Hu, Y. Miao, H. Ma, J. Gao, and J. Huang.
    Heterogeneous Aligned Fusion for Survival Prediction with Missing Modalities.
    In \textit{Medical Imaging with Deep Learning (MIDL)}, 2026.
    (To appear). \href{https://openreview.net/pdf?id=XjyK23sUkI}{\color{blue}Paper}

    \bibitem{Zheng2025}
    \textbf{Zheng Zheng}, X. Ni, and P. Hong.
    Multiple Abstraction Level Retrieve Augment Generation.
    \textit{arXiv:2501.16952}, 2025.\href{https://arxiv.org/pdf/2501.16952}{\color{blue} Paper}

    \bibitem{Zhang2023}
    Pengtao Zhang, \textbf{Zheng Zheng}, and Junlin Zhang.
    FiBiNet++: Reducing model size by low rank feature interaction layer for CTR prediction.
    In \textit{Proceedings of the 32nd ACM International Conference on Information and Knowledge Management},
    CIKM' 23,
    pages 4425–4429, New York, NY, USA, 2023.
    Association for Computing Machinery.
    \href{https://dl.acm.org/doi/10.1145/3583780.3615242}{\color{blue} Paper}

    \end{thebibliography}
    \endgroup
        
\vspace{-7pt}

%-----------EXPERIENCE-----------
% \section{Research Experience}
\section{Research Experience}
    
  \resumeSubHeadingListStart
    \resumeSubheading
      {Multiple Abstraction Level Retrieval-Augmented Generation}{Apr. 2024 - Apr.2025}
      {Visiting Research Scientist with Prof. Pengyu Hong}{Waltham, MA}
      % \resumeItemListStart
        % \resumeItem{Designed and implemented an end-to-end \textbf{Retrieval-Augmented Generation (RAG)} system pipeline, encompassing data gathering, preprocessing, extraction, model deployment, and performance evaluation for optimization.}
        % \resumeItem{Collected and curated a dataset of over 8,000 documents using advanced web scraping techniques (\textbf{Requests, BeautifulSoup4, Selenium-wire}) and leveraged \textbf{Grobid} for parsing raw documents into structured, hierarchical nodes across four abstract levels using \textbf{Llama-index}.}
        % \resumeItem{Employed a \textbf{semantic} approach with \textbf{Linq-Embed-Mistral} to split long texts and utilized \textbf{MapReduce} with the \textbf{Vicuna-13B} model to distill them into concise key information, enhancing text processing efficiency.}
        % \resumeItem{Created specialized QAR (question-answer-reason) datasets with \textbf{GPT-4-mini} based on specific node contexts selected by the \textbf{Bert classification model} and artificial rules.}
        % \resumeItem{Utilized \textbf{Linq-Embed-Mistral} to create embeddings and performed dimensionality reduction and clustering analysis.}
        % \resumeItem{Designed a robust evaluation framework for the RAG system using Ragas metrics.}
        % \resumeItem{Improved RAGAS-based answer correctness by \textbf{25.73\%} over a standard RAG baseline.}
    \resumeItemListStart
    \resumeItem{Architected a hierarchical Retrieval-Augmented Generation (RAG) system with multi-level document abstraction to support coarse-to-fine retrieval.}
    
    \resumeItem{Built a large-scale document ingestion and parsing pipeline (8,000+ PDFs) using \textbf{Requests}, \textbf{BeautifulSoup4}, \textbf{Selenium-wire}, and \textbf{Grobid}, converting raw documents into structured nodes across four abstraction levels via \textbf{LlamaIndex}.}
    
    \resumeItem{Designed semantic segmentation and embedding workflows with \textbf{Linq-Embed-Mistral}, and implemented MapReduce-based summarization using \textbf{Vicuna-13B} to compress long-context information.}
    
    \resumeItem{Constructed QAR (question-answer-reason) datasets using \textbf{GPT-4-mini} with rule-based and BERT-assisted node selection strategies to enhance supervision quality.}
    
    \resumeItem{Developed evaluation and benchmarking framework using \textbf{RAGAS} metrics, achieving a \textbf{25.73\%} improvement in answer correctness over standard RAG baselines.}
    
    \resumeItemListEnd
    \vspace{-2pt}
    
    \resumeSubheading
      {FiBiNet++ - Recommended System}{2022 - 2023}
      {Research Fellow}{Remote}
      \resumeItemListStart
        \resumeItem{ Conducted research focused on optimizing the model structure of both \textbf{bi-linear} and \textbf{SENet} structures within Recommandation System model \textbf{FibiNet} based on \textbf{Tensorflow} framework}
        \resumeItem{ Designed and implemented an innovative "Low Rank Layer" and also transformed the \textbf{element-wise Hadamard product} into an \textbf{inner product} method within the bi-linear interaction, leading to a significant reduction in the size of the FiBiNet model}
        \resumeItem{ Enhanced the efficiency of the "squeeze" step in the SENet structure by strategically splitting data points into multiple groups, improving the model's ability to capture one-dimensional features }
        \resumeItem{ Conducted rigorous benchmarking and performance evaluations of various models including \textbf{DNN}, \textbf{DeepFM}, \textbf{xDeepFM}, \textbf{DCN}, \textbf{AutoInt}, \textbf{DCN v2}, \textbf{FiBiNet}, and the newly developed \textbf{FiBiNet++}, achieving significant improvements in training and inference efficiencies, ranging from \textbf{37.50\% to 81.03\%} across the datasets}
        \resumeItem{Achieved a substantial reduction in model size for \textbf{FiBiNet}, decreasing it by \textbf{12x to 16x} across datasets, while also delivering remarkable performance improvements over state-of-the-art models}
    \resumeItemListEnd
    \resumeSubHeadingListEnd
    \vspace{-5pt}
\vspace{-7pt}

%-----------EXPERIENCE-----------
% \section{Research Experience}
\section{Teaching Experience}
    
  \resumeSubHeadingListStart
    \resumeSubheading
      {University of Texas at Arlington}{Jan. 2026 - Present}
      {Teaching Assistant}{Arlington, TX}
      \resumeItemListStart
        \resumeItem{ CSE-1325 : Object-Oriented Programming (Java) - Spring 2026 }
      \resumeItemListEnd
    \resumeSubheading
      {Brandeis Computer Science Department}{Sep. 2022 - May. 2023}
      {Teaching Assistant with Prof. Timothy J. Hickey}{Waltham, MA}
      \resumeItemListStart
        \resumeItem{ COSI-10A : Introduction to Problem Solving in Python - Fall 2022 \& Summer 2023}
        \resumeItem{ COSI-153A : Mobile Application Development - Summer 2023} 
        \resumeItem{ COSI-29A : Discrete Structures - Fall 2023 } 
      \resumeItemListEnd
    \resumeSubHeadingListEnd
\vspace{-7pt}

%-----------PROJECTS-----------
\section{Projects}
    \resumeSubHeadingListStart
        \resumeProjectHeading
          {\textbf{Hire Me Now} $|$ \emph{Docker, AWS, MERN stack, Django, Nginx, Certbot, Stripe, RESTful, LLM model}}{Mar. 2023}
          \vspace{-6pt}
          \resumeItemListStart
            \resumeItem{Designed and developed a full-stack website powered by \textbf{LLM} models aimed at assisting employees in their  job search competition, utilizing \textbf{Django} framework, and \textbf{MERN} stack (MongoDB, Express.js, React, Node.js)}
            \resumeItem{Built \textbf{Express} and \textbf{Django} microservices with serverless \textbf{REST API}. The Django implemented with \textbf{PEP8} code style}
            \resumeItem{Architected a React frontend with \textbf{Redux}, \textbf{Google Authentication} and \textbf{Stripe} API to forbid 100\% unauthorized access and enhance data status management}
            \resumeItem{Configured \textbf{Nginx} and \textbf{Certbot} with \textbf{SSL-certification} to facilitate communications between different services}
            \resumeItem{Containerized the entire system using \textbf{Docker} and developed \textbf{shell scripts} to automate installation and deployment}
            \resumeItem{Deployed the whole system on the \textbf{AWS EC2} service and incorporated a domain name using \textbf{GoDaddy} and \textbf{Route 53}}
          \resumeItemListEnd
        \vspace{-5pt}
      \resumeProjectHeading
          {\textbf{Multimodal QA system} $|$ \emph{LLM, Transformer, CNN, Decoder}}{Nov. 2024}
          \vspace{-6pt}
          \resumeItemListStart
            \resumeItem{Reproduced and analyzed the \textbf{multimodal vision-language model}, exploring its architecture and understanding the integration of visual and textual components for vision-language tasks.}
            \resumeItem{Fine-tuned the model on question answering benchmark \textbf{VQAv2} and \textbf{ScienceQA}, achieving state-of-the-art results across various benchmarks.}
            \resumeItem{Implemented and tested the multimodal pre-training process for PaliGemma, experimenting with different image resolutions (224px, 448px, and 896px) to optimize performance for high-resolution image tasks.}
            
          \resumeItemListEnd
        \vspace{-5pt}
         
    \resumeSubHeadingListEnd
    \resumeSubHeadingListStart
      \resumeProjectHeading
          {\textbf{Predicting Effective Arguments} $|$ \emph{RNN, GRU, LSTM, Bert, Pytorch, Matplotlib}}{Mar. 2023}
          \vspace{-6pt}
          \resumeItemListStart
            \resumeItem{Preprocess datasets into Torchtext data iterators, including steps such as \textbf{lowercasing}, \textbf{removing stopwords}, performing \textbf{stemming} and \textbf{lemmatization}, and handling \textbf{out-of-vocabulary (OOV)} words.}
            \resumeItem{Designed and trained various neural network models, including \textbf{RNN}, \textbf{LSTM}, \textbf{GRU}, and \textbf{BERT}, to classify argumentative elements within essays written by U.S. students.}
            \resumeItem{Conducted an ablation study combined with a grid search to systematically evaluate model performance, achieving the lowest loss and highest accuracy across different architectures.}
            
          \resumeItemListEnd
        \vspace{-5pt}
         
    \resumeSubHeadingListEnd
\vspace{-7pt}

%
%-----------PROGRAMMING SKILLS-----------
\section{Technical Skills}
\begin{itemize}[leftmargin=0.15in, label={}]
  \item \small{
    \textbf{Programming:} Python (advanced), Java, SQL, Shell, Go \\
    \textbf{Deep Learning:} PyTorch, TensorFlow, Distributed Training (DDP), Mixed Precision (AMP), Gradient Optimization, Multi-Task Learning, FlashAttention \\
    \textbf{LLM Engineering:} Transformers (HuggingFace), Fine-tuning (LoRA/PEFT), RAG Systems, Embedding Models, Prompt Optimization, Evaluation Frameworks, Multimodal Systems \\
    \textbf{Data \& Modeling:} Large-scale data preprocessing, Feature Engineering, Vector Indexing, Clustering \\
    \textbf{Systems \& Infrastructure:} Docker, AWS (EC2), Linux, GPU Training, Experiment Tracking, Git
  }
\end{itemize}

\end{document}
